% =============================================================
%  Chapter 6 — Discussion
% =============================================================
%  Interprets Chapter~\ref{ch:results} rather than re-deriving it:
%  no new numbers should be introduced here that are not already
%  in the Evaluation chapter.
% =============================================================

\chapter{Discussion}
\label{ch:discussion}


% ----- §6.1 Interpretation -----
\section{Interpreting the Results}
\label{sec:discussion-interpretation}

\paragraph{A methodological lesson that turned out to matter more than
any single capability finding.} The clearest pattern in
Chapter~\ref{ch:results} is not, in the end, a difference between models
at all: every cell that never reached a working baseline in this dataset
--- and most of the cells with only partial trajectories --- traces to
provider billing or quota exhaustion during data collection, not to
generation quality (Section~\ref{sec:results-rq3}). This is worth
dwelling on briefly, because the raw completion funnel
(Figure~\ref{fig:results-completion-funnel}) \emph{looks}, at a glance,
exactly like the kind of capability gap this thesis expected to find, and
an earlier draft of this chapter read it that way before the underlying
\texttt{llm\_error} text was actually checked cell by cell. The lesson
generalises beyond this dataset: a coarse failure-kind label
(\texttt{llm\_call}) is not a reliable proxy for model reliability unless
its error text is inspected, because it silently conflates transport
faults, provider-side account state, and (separately again) the
framework's own configured cost ceiling
(Section~\ref{sec:results-rq6}). None of those three is a statement about
the model.

\paragraph{What can and cannot be said about cost.} \texttt{glm-5.2} is
still roughly an order of magnitude cheaper per cell than
\texttt{gpt-5.5} or \texttt{claude-opus-4-8} (Section~\ref{sec:results-rq3}),
but that figure is now known to be inflated in \texttt{glm-5.2}'s and
\texttt{claude-opus-4-8}'s favour by cells that were never really
attempted: a rejected billing call costs almost nothing, so a model with
more billing-interrupted cells looks artificially cheap per cell. What
survives this caveat is narrower: on the cells every model was genuinely
able to attempt, all three reach a working baseline and go on to refine
it by a comparable multiplicative factor (Section~\ref{sec:results-rq1}).
Whether \texttt{glm-5.2}'s cost advantage comes with a real reliability
cost, a real reliability advantage, or no difference at all is exactly
the question the confounded cells need to be re-run to answer.

\paragraph{Deployment-specification versus code refinement.} The
division of labour in Section~\ref{sec:results-rq2} --- code refinement
carrying correctness-adjacent fixes with a heavy positive tail,
deployment refinement acting as a smaller and occasionally negative
capacity knob --- is consistent with what each lever can structurally
see. A spec-refinement step can move replicas, resources, and database
topology, but it cannot change what the running application \emph{does};
if the bottleneck is a connection-handling bug (Section~\ref{sec:results-rq6}),
no amount of deployment tuning reaches it. The reverse is not quite
symmetric: code refinement \emph{can}, in principle, work around a
deployment problem by changing application behaviour (for example,
reducing per-request work), which may explain why code refinement's
observed effect is not just larger on average but more variable in both
directions than deployment refinement's.

\paragraph{Convergence patterns.} Where a trajectory has enough successful
iterations to work with, Section~\ref{sec:results-rq5}'s
\textsc{ClickCount} case shows deployment parameters moving in a
consistent, workload-appropriate direction across several iterations
rather than oscillating --- once the agent identifies that horizontal
fan-out beats per-pod headroom for a near-stateless workload, later spec
iterations continue in that direction rather than re-deriving it from
scratch, which is the behaviour the append-only conversational memory of
Section~\ref{sec:methods-workflow} is designed to support. The backfire
case in the same section is the counter-example: convergence toward a
\emph{locally} reasonable decision is not the same as convergence toward
a good one, particularly when one refinement step is free to change
several spec fields whose interaction it cannot fully anticipate.


% ----- §6.2 Limitations -----
\section{Limitations}
\label{sec:discussion-limitations}

\paragraph{Scenario coverage.} Four of \textsc{BaxBench}'s 28 scenarios
were evaluated (Section~\ref{sec:setup-scenarios}), chosen to span a range
of workload shapes rather than to be a representative sample in a
statistical sense. All four require a database; none is compute-bound in
the way an image-processing or cryptographic scenario would be, so the
deployment-tuning story this thesis tells is specifically a
\emph{database-and-network-bound} story. Whether the same code/deployment
division of labour holds for CPU-bound workloads is untested.

\paragraph{Iteration budget.} Every experiment used a fixed $N{=}10$
refinement budget (Section~\ref{sec:setup-procedure}). Several
trajectories in Chapter~\ref{ch:results} were still improving, or had
just recovered from a regression, at iteration 010; a larger budget could
change both the achieved goodput ceiling and the RQ2 code-vs-spec split,
since which lever a decision stage reaches for next may depend on how
many refinement rounds it has already spent recovering from earlier
failures.

\paragraph{Sample count.} Every cell is $n=1$ (Section~\ref{sec:setup-procedure}).
A single code-generation seed can be unusually lucky or unlucky, and
Chapter~\ref{ch:results} already flags several ratios (Section~\ref{sec:results-rq1})
that rest on two-to-eight healthy-baseline cells per model. Every
comparative number in this thesis should be read as a single-draw
estimate, not a distribution; nothing here has been significance-tested
against sampling variance, because there is no repeated-sample data to
test it against.

\paragraph{Benchmark methodology.} Sustained goodput
(Section~\ref{sec:methods-load-profile}) is a single scalar chosen to make
iterations comparable without per-scenario rate tuning. It necessarily
discards information a multi-metric evaluation would keep --- tail
latency distribution shape, resource-utilisation efficiency, cost per
served request --- some of which is already collected as diagnostics
(Section~\ref{sec:methods-outcome}) but not used as a scoring signal.
The explore/recovery/refine load-profile mechanics
(Section~\ref{sec:methods-load-profile}) also mean that a deployment which
never passes its warm-up health check is scored identically (goodput
$=0$) to one that technically starts but collapses immediately under the
first load step; Section~\ref{sec:results-rq1} treats both as ``dead
baselines'' without distinguishing them further.

\paragraph{Stochastic LLM behaviour.} Temperature was fixed at $0.2$
(Section~\ref{sec:setup-models}), a deliberate reduction of but not an
elimination of run-to-run variance; combined with $n=1$ sampling, this
thesis cannot separate ``this model is worse at this scenario'' from
``this particular draw was unlucky.'' Provider-side model updates
between the writing and the final submission of this thesis are a
related, less controllable source of drift: model snapshots were pinned
where providers exposed a dated snapshot (Section~\ref{sec:setup-models}),
but not every provider does so.

\paragraph{Data-collection interruptions.} Thirteen of the 36 cells in
this dataset were touched by an LLM-provider billing or quota exhaustion
event during data collection, and ten never reached a working baseline as
a direct result (Section~\ref{sec:results-rq3}). This is not a property
of the framework or the models under test; it is an operational
limitation of running a compute- and API-budget-intensive evaluation on
finite accounts, and it currently makes this thesis's model-comparison
and framework-comparison claims (RQ3, RQ4) substantially weaker than a
first read of the completion funnel would suggest. It is listed here as
its own limitation, separate from stochastic LLM behaviour above, because
it is not sampling noise: it is missing data that a re-run with a
sufficiently funded account would recover deterministically, not average
away.

\paragraph{Infrastructure constraints.} All results come from one cluster
topology and one hardware profile (Section~\ref{sec:setup-infrastructure}).
``Fixed infrastructure constraints'' in the thesis research question
(Section~\ref{sec:intro-missing}) is therefore fixed to \emph{one}
specific capacity envelope; whether the same deployment decisions
(e.g.\ the replica-count/CPU-request trade-off of Section~\ref{sec:results-rq5})
would be reachable, or would be the right decisions, on a cluster an
order of magnitude larger or smaller is untested.


% ----- §6.3 Practical Implications -----
\section{Practical Implications}
\label{sec:discussion-implications}

\paragraph{For agentic software optimisation.} The clearest transferable
result is the code/deployment division of labour itself
(Section~\ref{sec:results-rq2}): when an agent can act on two levers that
address different failure classes, keeping them mutually exclusive per
step (Section~\ref{sec:methods-workflow}) is what makes it possible to
attribute a measured change to a cause at all, and this thesis's data
suggests the two levers really do specialise in practice, not just by
construction. Systems that let an agent freely mix implementation and
configuration changes in one step should expect the attribution problem
Section~\ref{sec:results-rq6}'s backfire case illustrates \emph{within} a
single lever, only worse.

\paragraph{For Kubernetes deployment optimisation.} The
\textsc{ClickCount} trajectory (Section~\ref{sec:results-rq5}) is a small
existence proof that an LLM agent, given only a schema and raw load-test
feedback, can rediscover a standard scaling heuristic (favour horizontal
fan-out over per-pod headroom for CPU-light, network-bound workloads)
without being told it. That is a meaningfully lower bar than fully
autonomous production capacity planning, but it is a higher bar than
today's reactive autoscalers (Section~\ref{sec:rw-cluster}) clear, since
HPA/VPA-style mechanisms do not reason about database topology or
placement at all.

\paragraph{For LLM-guided performance tuning more broadly.} The
recurring finding that correctness bugs dominate the failure budget
(Section~\ref{sec:results-synthesis}) is a caution against evaluating
performance-tuning agents purely on their tuning behaviour in isolation:
in this dataset, the single largest lever on final goodput was not a
tuning decision at all, but whether the code was correct enough to
deploy under load in the first place --- a class of bug that, as
Section~\ref{sec:results-rq6}'s pooling case study shows, a
single-container functional-test suite cannot even see.
